% Book-local P07.1 obstruction, 20 September 2026.
% This isolated statement is not yet a live-book input and makes no novelty
% or priority claim.

\begin{proposition}
\label{prop:p071-trivalent-double-junction-obstruction}
Let $S$ and $T$ be node neighbourhoods in two exact trivalent
Chambert--Loir--Ducros stars.  Use outward coordinates
$r,s\in[0,R)$ on their branches, label both branch sets by
$\{1,2,3\}$, and write $o=(\eta_S,\eta_T)$.  On $S\times T$, let
$\mathcal G_0$ be the following local two-star generator sheaf.  A member
$g$ is real-valued, has smooth jointly convex restrictions $g_{ij}$ to the
nine branch rectangles, has common smooth restrictions on the two families
of node axes, and satisfies
\begin{align}
 \sum_{i=1}^3\partial_rg_{ij}(0,s)&=0
       &&(j\in\{1,2,3\},\ 0\leq s<R),
 \label{eq:p071-trivalent-first-balance}\\
 \sum_{j=1}^3\partial_sg_{ij}(r,0)&=0
       &&(i\in\{1,2,3\},\ 0\leq r<R).
 \label{eq:p071-trivalent-second-balance}
\end{align}
Let $\mathcal P$ be the smallest subsheaf containing $\mathcal G_0$ and
closed under local pointwise decreasing sequential limits.

Every $v\in\mathcal P$ satisfies the following mixed double-junction
inequality near $o$.  Take finitely many atoms
\[
 z_a=(p_{i_a}(r_a),q_{j_a}(s_a)),
 \qquad \lambda_a>0,
 \qquad \sum_a\lambda_a=1,
\]
in one product node chart.  If there are $M,N\geq0$ such that
\begin{align}
 \sum_{a:i_a=i}\lambda_ar_a&=M
       &&(i=1,2,3),
 \label{eq:p071-trivalent-first-moment}\\
 \sum_{a:j_a=j}\lambda_as_a&=N
       &&(j=1,2,3),
 \label{eq:p071-trivalent-second-moment}
\end{align}
then
\begin{equation}
                       v(o)\leq\sum_a\lambda_av(z_a).
 \label{eq:p071-trivalent-mixed-jensen}
\end{equation}
The assertion uses the usual extended-real convention and is automatic
when $v(o)=-\infty$.

This mixed condition is not implied by joint convexity on the nine branch
rectangles together with both one-factor star--line junction families.
Indeed, on $[0,R)^2$ define
\begin{equation}
 u_{ij}(r,s)=r+s-3\delta_{ij}\min\{r,s\}.
 \label{eq:p071-trivalent-obstruction-function}
\end{equation}
The nine functions glue to a continuous real-valued function $u$ on
$S\times T$.  Every $u_{ij}$ is jointly convex.  For each fixed second
branch $j$, $u$ satisfies every first-star--line junction inequality:
whenever $\lambda_a>0$, $\sum_a\lambda_a=1$,
\begin{equation}
 \sum_{a:i_a=i}\lambda_ar_a=M\quad(i=1,2,3),
 \qquad s_0=\sum_a\lambda_as_a,
 \label{eq:p071-trivalent-first-separate-data}
\end{equation}
with all atoms in that branch, one has
\begin{equation}
 u(\eta_S,q_j(s_0))
       \leq\sum_a\lambda_a
             u(p_{i_a}(r_a),q_j(s_a)).
 \label{eq:p071-trivalent-first-separate-jensen}
\end{equation}
The symmetric second-star--line family also holds.  Nevertheless, taking
the three atoms $(p_i(\rho),q_i(\rho))$, each with weight $1/3$, for any
$0<\rho<R$, gives
\begin{equation}
 u(o)=0> -\rho
   =\frac13\sum_{i=1}^3u(p_i(\rho),q_i(\rho)).
 \label{eq:p071-trivalent-obstruction-violation}
\end{equation}
Consequently $u\notin\mathcal P$ near $o$.  Thus any comparison class for
the genuinely branched two-star case must include mixed double-junction
conditions; iterating the two one-star comparison tests is insufficient.
This proposition does not assert that the mixed conditions are sufficient,
and it does not retire P07.1.
\end{proposition}
